% Power-balance composition at T = 6, case 6e (all bounds imposed).
%
% eq:cp_balance says P_Subs^t + P_B^t - p_L^t = 0. Splitting the battery by sign,
%   supply side   P_Subs^t + P_d^t,   P_d = max(P_B, 0)   (discharging)
%   demand side   p_L^t    + P_c^t,   P_c = max(-P_B, 0)  (charging)
% and the two sums are equal at every interval. The figure draws them as a pair
% of stacked bars per interval: equal total height is the balance constraint, and
% the segment split is each asset's share of the kW mix.
%
% Data comes from figures/balance.csv, written by make_balance.jl out of the
% VERIFIED centralized reference (results/copper_plate/centralized_reference.csv,
% case 6e_T6), not hand-copied. The generator asserts the balance to 1e-7 before
% writing, so the figure's central claim is checked in code, not by eye.
%
% Three axes are overlaid because pgfplots stacks every ybar in one axis
% together: axis 1 carries the supply stack and all ticks, axis 2 the demand
% stack, axis 3 the equal-height rule. All three share identical limits so the
% bars register.
%
\definecolor{cSubsB}{HTML}{31708E}  % steel blue -- substation import
\definecolor{cDisB}{HTML}{800000}   % maroon     -- discharging (as in schedule_fig)
\definecolor{cLoadB}{HTML}{C08552}  % ochre      -- served demand
\definecolor{cChgB}{HTML}{008000}   % green      -- charging    (as in schedule_fig)
\definecolor{cRule}{HTML}{1A1A1A}   % near-black -- equal-height rule
\definecolor{gridMajorB}{HTML}{E8C4D4}
\definecolor{gridMinorB}{HTML}{D4D4E8}
%
\begin{figure}[!t]
\centering
\begin{tikzpicture}
\pgfplotsset{
  balanceaxis/.style={
    width=0.80\columnwidth, height=0.46\columnwidth, scale only axis,
    xmin=0.45, xmax=6.55, ymin=0, ymax=2450,
    xtick={1,2,3,4,5,6},
    ytick={0,500,1000,1500,2000},
    tick align=outside,
    tick label style={font=\scriptsize},
    label style={font=\scriptsize},
    axis line style={black, line width=0.6pt},
  },
}

% ---------------- axis 1: supply side (P_Subs + P_d), carries the ticks -------
\begin{axis}[
    balanceaxis, name=axBal,
    ybar stacked, bar width=13pt,
    grid=both,
    major grid style={line width=0.6pt, color=gridMajorB, opacity=0.7},
    minor grid style={line width=0.4pt, color=gridMinorB, opacity=0.6},
    xlabel={Time interval $t$},
    ylabel={Power [kW]}, ylabel near ticks,
    legend style={at={(0.5,1.04)}, anchor=south, legend columns=4,
                  draw=none, fill=none, column sep=0.7ex, font=\scriptsize},
    legend image code/.code={\draw[#1] (0cm,-0.085cm) rectangle (0.24cm,0.085cm);},
    legend cell align=left,
]
\addplot[draw=cSubsB!70!black, fill=cSubsB, fill opacity=0.85, line width=0.5pt]
  table[col sep=comma, x=tsup, y=psub] {figures/balance.csv};
\addlegendentry{$P_{\text{Subs}}$}
\addplot[draw=cDisB!70!black, fill=cDisB, fill opacity=0.85, line width=0.5pt]
  table[col sep=comma, x=tsup, y=pd] {figures/balance.csv};
\addlegendentry{$P_B > 0$ (discharge)}
\addlegendimage{draw=cLoadB!70!black, fill=cLoadB, fill opacity=0.85}
\addlegendentry{$p_L$}
\addlegendimage{draw=cChgB!70!black, fill=cChgB, fill opacity=0.85}
\addlegendentry{$P_B < 0$ (charge)}
\end{axis}

% ---------------- axis 2: demand side (p_L + P_c), overlaid ------------------
\begin{axis}[
    balanceaxis, at={(axBal.south west)}, anchor=south west,
    ybar stacked, bar width=13pt,
    axis lines=none, xtick=\empty, ytick=\empty, grid=none,
]
\addplot[draw=cLoadB!70!black, fill=cLoadB, fill opacity=0.85, line width=0.5pt]
  table[col sep=comma, x=tdem, y=pl] {figures/balance.csv};
\addplot[draw=cChgB!70!black, fill=cChgB, fill opacity=0.85, line width=0.5pt]
  table[col sep=comma, x=tdem, y=pc] {figures/balance.csv};
\end{axis}

% ---------------- axis 3: the equal-height rule ------------------------------
\begin{axis}[
    balanceaxis, at={(axBal.south west)}, anchor=south west,
    axis lines=none, xtick=\empty, ytick=\empty, grid=none,
]
\addplot[only marks, mark=-, mark size=15.5pt, cRule, line width=0.9pt]
  table[col sep=comma, x=t, y=total] {figures/balance.csv};
\end{axis}
\end{tikzpicture}
\caption{Where every kW comes from and goes, at $T = 6$ with all bounds imposed.
Each interval carries two stacked bars: on the left what \emph{supplies} the
node, substation import $P_{\text{Subs}}^t$ plus battery discharge; on the right
what \emph{consumes} it, the served demand $p_L^t$ plus battery charging.  The
two must reach the same height --- that equality \emph{is} the power balance
\eqref{eq:cp_balance}, and the black rule marks the level they share.  Read the
figure as one worked example.  In $t = 2$ the price peaks at $0.197$\,\$/kWh, and
the battery discharges at its $450$\,kW rail, so the maroon block supplies
$23.8\%$ of the $1888$\,kW the node needs and the substation is held down to
$1438$\,kW even though demand is near its own peak.  Three intervals later the
sign flips: at $t = 5$ the price bottoms at $0.083$\,\$/kWh, the battery charges
at $-450$\,kW, and that charging appears on the \emph{demand} side, pushing
substation import up to $2073$\,kW to serve only $1623$\,kW of actual load.  The
device is doing arbitrage in the open.  Over the horizon the substation imports
$10\,606$\,kWh against $10\,657$\,kWh of load served, the $51$\,kWh difference
being the net the battery gives back as it returns from $B = 1200$\,kWh to
$1950$\,kWh.  Bars are in kW on the tADMM base, $1$\,p.u.\ $= 1000$\,kW.}
\label{fig:balance}
\end{figure}
